\section{Limitations and Possible Improvements}

The submitted desktop MVP is intentionally limited to two-dimensional Static Analysis and Modal Analysis. It is an educational tool rather than a commercial structural analysis package. It does not claim three-dimensional modeling, nonlinear analysis, design-code checking, staged construction, automated meshing, advanced load combinations, production-quality drawing tools, or final desktop RSA/THA workflows.

Main limitations are the two-dimensional DOF set \([u_x,u_y,r_z]\), the educational lumped-mass-oriented modal workflow, the deliberately transparent standard-library numerical core, and remaining placeholder figures/tables that must be replaced before final packaging. More advanced modelling features were either not implemented or only partly explored because they were not necessary for the submitted Static and Modal scope. These include translational and rotational spring supports, member axis transformation options, inclined supports, rigid end zones, three-dimensional model generation, torsional deformation, and full three-dimensional member transformation.

Possible improvements include a polished installer, richer classroom examples, automated report generation, additional benchmark cases, expanded plot styling, more robust input checking, and broader modal modeling options. The omitted or partial modelling features could also be implemented in future work to investigate different structural cases, especially models with support flexibility, skewed boundary conditions, member end offsets, torsion-sensitive behaviour, or three-dimensional framing action. RSA and THA should be promoted to desktop workflows only if the project scope is formally expanded, and any future feature should preserve the same model-to-assembly-to-solver-to-result-object-to-visualization architecture.
